\documentclass[main.tex]{subfiles}


\begin{document}

%from pagenumber to up
% \phantomsection 
% \hypertarget{toc}{}

 

 

 
   

\section{Диэлектрик}


\textbf{Диэлектрик}~--- это тело, \emph{не} способное проводить через себя электрические заряды. Например, если металлическое \emph{заряженное} тело~А соединить диэлектриком с металлическим \emph{незаряженным} телом~Б, то заряд тела~А \emph{не} перераспределится между этими двумя телами.

В диэлектрике (в отличие от проводника) \emph{нет свободных зарядов}. Заряженные частицы (электроны и ядра) в диэлектрике являются \emph{связанными}~--- электроны при этом могут перемещаться лишь внутри молекулы тела (рис.\,\ref{fig:p113}).

    \begin{figure}[H]
    \centering

    \begin{tikzpicture}[font=\small,            line cap=round,                             >={Stealth[inset=0pt,round,length=3mm, angle'=20]},vec/.style={>={Stealth[inset=0pt,round,length=3.5mm, angle'=20]}}]


\tikzset{atom/.pic={
  \begin{scope}[shift={({rnd*360}:.05)}]
  \shade[ball color=red] (0,0) circle (\dn);
  \shade[ball color=NavyBlue!70,rotate={rnd*360}] ([shift={(180:{\dn cm})}]0,0) circle (\de);
  \end{scope}}
}

\tikzset{atompo/.pic={
  \begin{scope}[shift={({rnd*360}:.05)}]
  \shade[ball color=red] (0,0) circle (\dn);
  \end{scope}}
}

\tikzset{atomne/.pic={
  \begin{scope}[shift={({rnd*360}:.05)},rotate={rnd*360}]
  \shade[ball color=red] (0,0) circle (\dn);
  \shade[ball color=NavyBlue!70] ([shift={(180:{\dn cm})}]0,0) circle (\de);
  \shade[ball color=NavyBlue!70] ([shift={(0:{\dn cm})}]0,0) circle (\de);
  \end{scope}}
}

\def\dn{.15}
\def\de{.1}


%solid
\fill[yellow,nearly transparent] (0,0) rectangle++ (4,2);
\draw[] (0,0) rectangle++ (4,2);


%atoms left
\foreach \y in {0,...,2}
\foreach \x in {0,...,5}
\pic[shift={(.4,.4)}] at ({0+(4-.8)/5*\x},{0+(2-.8)/2*\y}) {atom};



    \end{tikzpicture}
\caption{Диэлектрик}
\label{fig:p113}
\end{figure}

Соответствующие электроны (синие шары) и ядра (красные шары) сцеплены друг с другом (если диэлектрик жидкий или газообразный, то его молекулы <<носят с собой>> свои ядра и электроны)\footnote{В особых случаях электрон может <<оторваться>> от своего ядра~--- так происходит, например, при электризации диэлектрика трением.}.

Пусть диэлектрик помещен в однородное электрическое поле (рис.\,\ref{fig:p114}).

\tikzfading[name=fade right,
left color=transparent!0,
right color=transparent!100]

\tikzfading[name=fade left,
left color=transparent!100,
right color=transparent!0]

    \begin{figure}[H]
    \centering

    \begin{tikzpicture}[font=\small,            line cap=round,                             >={Stealth[inset=0pt,round,length=3mm, angle'=20]},vec/.style={>={Stealth[inset=0pt,round,length=3.5mm, angle'=20]}}]


\tikzset{atom/.pic={
  \begin{scope}[shift={({rnd*360}:.05)}]
  \shade[ball color=red] (0.05,0) circle (\dn);
  \shade[ball color=NavyBlue!70,rotate={rand*30}] ([shift={(180:{\dn cm})}]0.05,0) circle (\de);
  \end{scope}}
}

\tikzset{atompo/.pic={
  \begin{scope}[shift={({rnd*360}:.05)}]
  \shade[ball color=red] (0,0) circle (\dn);
  \end{scope}}
}

\tikzset{atomne/.pic={
  \begin{scope}[shift={({rnd*360}:.05)},rotate={rnd*360}]
  \shade[ball color=red] (0,0) circle (\dn);
  \shade[ball color=NavyBlue!70] ([shift={(180:{\dn cm})}]0,0) circle (\de);
  \shade[ball color=NavyBlue!70] ([shift={(0:{\dn cm})}]0,0) circle (\de);
  \end{scope}}
}

\def\dn{.15}
\def\de{.1}


%solid
\fill[yellow,nearly transparent] (0,0) rectangle++ (4,2);
\fill[NavyBlue,path fading=fade right] (0,0) rectangle++ (.5,2);
\fill[red,path fading=fade left] (4,0) rectangle++ (-.5,2);
\draw[] (0,0) rectangle++ (4,2);


%atoms left
\foreach \y in {0,...,2}
\foreach \x in {0,...,5}
\pic[shift={(.4,.4)}] at ({0+(4-.8)/5*\x},{0+(2-.8)/2*\y}) {atom};

%E
\draw[->,Green] (-2,2.5) --++ (8,0) node [right,black] {$E_\text{в}$};
\draw[->,Green] (-2,-.5) --++ (8,0) node [right,black] {$E_\text{в}$};


%phantom
\phantom{\node[left] at (-2,2.5) {$E_\text{в}$};}
\phantom{\node[left] at (-2,-.5) {$E_\text{в}$};}


    \end{tikzpicture}
\caption{Диэлектрик в электрическом поле}
\label{fig:p114}
\end{figure}

Внешенее поле~$E_\text{в}$ <<разворачивает>> молекулы диэлектрика так, что на одной поверхности тела (левая поверхность~--- синяя область) оказываются преимущественно отрицательные заряды ~--- электроны, а на другой поверхности (правая поверхность~--- красная область) оказываются в основном положительные заряды~--- ядра. Это происходит вследствие действия силы Кулона на электроны со стороны поля: на рис.\,\ref{fig:p114} эта сила <<тянет>> электроны влево.  

Избыточные \emph{наведенные} заряды на поверхностях диэлектрика (см.~рис.\,\ref{fig:p114}) создают внутри тела собственное поле~$E_i$, направленное \emph{против} внешнего поля~$E_\text{в}$. Поле~$E_i$ ослабляет поле~$E_\text{в}$ внутри диэлектрика (при этом $E_i<E_\text{в}$). Результирующее (суммарное) поле~$E$ внутри тела равно $E=E_\text{в}-E_i>0$. 

\textbf{Диэлектрическая проницаемость} ($\varepsilon$)~--- это характеристика тела, показывающая во сколько раз уменьшается поле в нем по сравнению с вакуумом:
\begin{equation}
    \varepsilon=\frac{E_\text{в}}{E},
\end{equation}
где $E_\text{в}$~--- поле в вакууме, $E$~--- поле в данном теле ($E<E_\text{в}$).


 
\end{document}
